\section{Experimental Setup}
\label{sec:setup}

\paragraph{Tasks.}
Nine character tasks over English words of length 3--9: \texttt{first\_letter},
\texttt{last\_letter}, \texttt{contains\_letter} (yes/no), \texttt{extraction} ($k$-th letter),
\texttt{counting} (occurrences of a queried letter), \texttt{most\_frequent\_letter},
\texttt{letter\_before} (the letter preceding a queried letter), \texttt{word\_length} and
\texttt{count\_distinct}. Two further tasks, \texttt{counting\_lm} and \texttt{count\_distinct\_lm},
ask the same questions as the two count tasks but sample examples so that word-only shortcuts
carry no information about the answer (\S\ref{sec:nuisance}, Appendix~\ref{app:lm}). Five tasks
(\texttt{first\_letter}, \texttt{last\_letter}, \texttt{most\_frequent\_letter},
\texttt{word\_length}, \texttt{count\_distinct}) are pure functions of the word; the others carry a
query letter or position. Prompts and answer sets are in Appendix~\ref{app:prompts}.

\paragraph{Stimuli and models.}
Each task has 800 examples sampled from a 6000-word frequency-ranked English vocabulary built from
the Brown Corpus \citep{kucera1967brown}, realizing 405--757 unique words per task. We evaluate 16
open-weight models from seven families: Qwen3 (0.6/1.7/4/8B), Llama-3.x (1/3/8B), Gemma-2
(2/9B), OLMo-2 (1/7/13B), Phi-4 (3.8B-mini, 14B), OLMo-Hybrid-7B and OLMoE-1B-7B
\citep{qwen3report,llama3report,gemma2report,phi4report,olmo2report,olmoereport,olmohybridreport};
Table~\ref{tab:models} (Appendix~\ref{app:tables}) lists them.
Six are post-trained (the four Qwen3 releases and both Phi-4 models); we report base-only
statistics over the other ten where post-training could matter. All run in bf16 on one 48\,GB
GPU. Prompts are completion-style, at 0 or 3 demonstrations. At zero shots no arm contains a
demonstration, which makes it the clean contrast; most claims are reported at both counts.

\paragraph{Prompt arms and positions.}
For every example we run three prompts on the same word and label (Figure~\ref{fig:overview}a):
the \textbf{question}, the task prompt (e.g.\ \texttt{The last letter of the word "cat" is});
a \textbf{neutral} sentence, \texttt{The word "cat" appears in this sentence.}; and a
\textbf{continuation} carrier, \texttt{Consider the word "cat". The answer is}, which matches the
question's ``awaiting an answer'' final position without naming the property. Term tasks carry
their query letter or position into both carriers. We read hidden states at two positions: the
\textbf{answer position} (the final prompt token) in every arm, and the \textbf{word position} in
the question and neutral arms, for the pure-word tasks, as either the word's last token or the mean
over all tokens overlapping the word.

\paragraph{Probe and estimator.}
A standardized multinomial logistic regression ($\ell_2 = 10^{-3}$) is fit at each of 12 evenly
spaced layers (including the embedding layer and the last layer), with three-fold cross-validation
grouped by query word, so no query word appears in both training and test labels. Folds are
cold-started. The reported accuracy uses \emph{held-out layer selection}: for each test fold, the
layer is chosen on the other folds. With this estimator, the peak-over-layers accuracy exceeds the
held-out-selected accuracy by 0.001--0.019 per task, with no positive dependence on model depth,
and all representation numbers lie within $\pm 0.021$ of an absolute-stride, warm-started
re-estimate (Appendix~\ref{app:estimator}).

\paragraph{Baselines.}
The \textbf{majority} baseline predicts the training fold's modal answer. The \textbf{nuisance
floor} is the best of several predictors that see only the word: its length, its position-free
letter set, the query letter where the task has one, and their combination, fit on the same folds.
Predictors that are sufficient for the label (length for \texttt{word\_length}, the letter set for
\texttt{count\_distinct}) are excluded. We report accuracy \emph{above the floor} unless stated
otherwise.

\paragraph{Readout.}
The model's own accuracy is the argmax of its next-token logits over the answer-token set (26
letters, yes/no, or ten digits). At zero shots, completion prompts are a degenerate format for
readout: larger models continue the sentence rather than answer, so readout claims use three shots.

\paragraph{Statistics.}
Scaling associations are Pearson correlations or OLS slopes against $\log_{10}$ parameters, with
two-sided permutation $p$-values from 10{,}000 permutations of model size. Because each model
contributes every arm, a question$\times$scale interaction equals the slope of the per-model
difference between arms. We also report a within-family (fixed-effect) slope, a 95\% interval
from resampling families, and base-only correlations. Where a claim spans the four
prompt-sensitive tasks we apply Holm correction \citep{holm1979}. The 16 checkpoints are related
and not independent draws; these are cross-model associations, not population scaling laws.
